\section{Background}
\label{sec:background}

\subsection{Conformational B-cell epitope prediction}

A B-cell epitope is the part of an antigen that an antibody binds in an
antibody--antigen complex. Linear B-cell epitopes are contiguous in sequence;
conformational, or discontinuous, epitopes, which are the large majority of
what antibodies actually recognise, are assembled from residues that the folded
structure brings together but that may be far apart in the chain. Prediction is
therefore per-residue binary labelling, each residue being called an epitope
residue or not, and two properties of that task are awkward. Epitope residues
are a minority of all residues, and the label depends on structural context
that the sequence encodes only implicitly.

Methods have moved through three broad generations. Propensity-scale approaches
scored residues with physicochemical heuristics. Structure-based methods used
solved or modelled structures to compute relative solvent accessibility and
protrusion for each surface residue.
The current generation feeds representations from a pre-trained protein language
model into a supervised head, which is how BepiPred-3.0 \citep{clifford2022bepipred}
and structure-aware successors such as DiscoTope-3.0 \citep{hoie2024discotope}
are built. Recent additions include the SEMA 2.0 web platform
\citep{ivanisenko2024sema} and residue-level predictors built directly on ESM-2
projections \citep{israeli2024single}. Both sit close to the PLM family
evaluated here. \citet{gabellieri2026bcell} survey the field as it now
stands. We use BepiPred-3.0 as the reference predictor on BP3C50ID, where it is the
strongest openly runnable sequence-based baseline and the method a laboratory
would realistically be choosing between. On the IEDB-derived subset we do not
have BepiPred-3.0 predictions, so the reference there is our own local-only
predictor. The two fusion arms therefore test different pairs, and we keep
them separate throughout.

Evaluation in this area is dominated by the non-redundant BP3C50ID benchmark
distributed with BepiPred-3.0, whose external test split holds
\BepipredNumAntigens{} antigens (\BepipredResidues{} residues,
\BepipredPositives{} of them positive). That independent benchmark is a small
one, and its size is worth keeping in view when reading any comparison built on
it, including ours. We report the Matthews correlation coefficient (MCC)
alongside the threshold-independent ROC-AUC and average precision, because MCC
behaves sensibly under the class imbalance that per-residue epitope labelling
produces
\citep{chicco2020mcc}. Under that imbalance, accuracy and F1 can both be high
for a predictor that has learned little.

\subsection{Splitting inference between a local device and a remote service}

Running part of a model locally and part remotely is an old idea with a current
name. The systems literature treats it as a scheduling problem. The task is to
decide where each input is processed, given a local
model, a remote model, and a per-input signal about difficulty, so that some
combination of accuracy, latency, cost and bandwidth is optimised
\citep{matsubara2022split}. The signal is usually the
local model's own confidence. That makes the mechanism attractively simple, and
it also introduces its characteristic failure, since the decision is only as
good as the calibration of the model making it.

Whether offloading saves energy at all is a settled question in that
literature. The answer depends on the ratio of computation offloaded to the
overhead of offloading it. Below a threshold ratio, offloading costs more energy
than it saves \citep{kumar2010cloud}. The same effect has been observed
recently in model cascades, where a draft-and-verify decoding scheme consumed
more energy than running the large model alone despite reducing latency
\citep{dutta2026speculative}. Our contribution on this axis is not the
phenomenon but the location of the threshold for this task, and the observation
that per-residue epitope screening sits on the unfavourable side of it.

Two things are different in the setting we study. The first is that the
quantity a laboratory wants to minimise is often not bandwidth or cost but
disclosure, meaning how many of its sequences are transmitted at all.
Disclosure is a counting measure rather than a continuous one, and no
accuracy-oriented objective captures it. The second is that the per-antigen
compute is small enough that the routing decision is not free relative to the
work it schedules.
In large-model serving, a confidence check is negligible beside the inference it
gates. In per-residue epitope screening on short sequences, it is not. The
carbon results in Section~\ref{sec:routing-results} are a consequence of that
ratio.

\subsection{Accounting for energy and carbon}

Reported carbon figures for machine-learning systems are rarely measurements.
The common practice, which we follow and label as such, is to derive energy from
a device power figure and a measured runtime. That energy is then converted to
mass of carbon dioxide equivalent with a grid-intensity constant
\citep{strubell2019energy,lacoste2019quantifying,patterson2021carbon}. Our
energy model follows the formulation of Green Algorithms
\citep{lannelongue2021green}, which also supplies the amortised-hardware term.
That formulation has been applied to bioinformatics workloads specifically
\citep{grealey2022carbon}. We claim no novelty in the instrument. The
result is a model whose output is proportional to measured latency by
construction. The model is useful for comparing policies that run on the same
hardware, since the proportionality constant cancels; it is not a measurement
of what any particular laboratory's equipment would draw, and we are careful
not to present
it as one. Table~\ref{tab:provenance} records which of our reported quantities
are measured and which are modelled.

\subsection{Selection on the evaluation split}

When a method has a tuning parameter, an ensemble weight or a decision
threshold or a fusion coefficient, and that parameter is chosen by scoring a
parameter grid on the data the paper reports, the reported figure is an
over-optimistic estimate of the predictor's true error. Model selection bias of
this kind grows with the number of grid points tried and shrinks with the size
of the evaluation set, and it is not a new observation.

The mechanism itself is one line of algebra, and it is worth setting down
because it fixes what can and cannot be concluded from it. Write
\(\mathcal{W}_G\) for a grid of \(G\) candidate values of the tuning
parameter, \(\widehat{M}(w)\) for the metric of the configuration indexed by
\(w\) computed on the split the figure is reported on, and \(M(w)\) for the
quantity that \(\widehat{M}(w)\) estimates. A grid scan reports
\(\max_{w \in \mathcal{W}_G} \widehat{M}(w)\). For any grid point
\(w_0 \in \mathcal{W}_G\) fixed in advance of seeing that split, and for any
finer grid \(\mathcal{W}_{G'} \supseteq \mathcal{W}_G\),
\begin{equation}
\label{eq:selection-max}
\mathbb{E}\Big[\max_{w \in \mathcal{W}_G} \widehat{M}(w)\Big]
  \;\ge\; \mathbb{E}\big[\widehat{M}(w_0)\big] \;=\; M(w_0),
\qquad
\mathbb{E}\Big[\max_{w \in \mathcal{W}_G} \widehat{M}(w)\Big]
  \;\le\; \mathbb{E}\Big[\max_{w \in \mathcal{W}_{G'}} \widehat{M}(w)\Big].
\end{equation}
Both parts hold split by split before any expectation is taken: the maximum of
a finite set is at least as large as any member of it, which gives the first
inequality, and enlarging the set over which a maximum is taken cannot decrease
it, which gives the second; expectation preserves both orderings. The single
equality is the one assumption, namely that at a grid point chosen without
reference to the split, \(\widehat{M}(w_0)\) is unbiased for \(M(w_0)\),
which is what a held-out split buys and what the maximum destroys.

Equation~\eqref{eq:selection-max} settles two things and no more. It fixes the
sign: a scanned maximum sits on average at or above the performance of any
configuration one could have committed to beforehand, so the direction of the
distortion is known before any data are seen. And it makes the effect
non-decreasing under refinement of a \emph{nested} family of grids, which is
why a wider grid is the more flattering protocol. That condition is worth
stating, because equally spaced grids of different sizes are not nested: the
weight \(0.25\) lies in a five-point grid but not in an eleven-point one, so
two such grids are not ordered by inclusion and the inequality says nothing
about that pair. It supplies no magnitude either. How far above \(M(w_0)\) the
expected maximum sits depends on the correlation and the noise scale of the
\(\widehat{M}(w)\) across the grid, which the inequality does not constrain;
the bound is not tight, and it holds with equality only where the scan returns
\(w_0\), or a grid point with the same \(\widehat{M}\), on every split. The magnitude reported in this paper is
therefore measured and not derived: Section~\ref{sec:bias-curve} estimates it by resampling, and
that measured gap is an estimate of the difference between two protocols on
these data, not of the expected bias in
Equation~\eqref{eq:selection-max}.

\citet{varma2006bias} demonstrated it for bioinformatics classifiers two decades
ago and showed that nested resampling, which estimates the error in an outer
loop over the left-out samples, returns a nearly unbiased estimate.
\citet{boulesteix2009overoptimism}
named the resulting pattern over-optimism. \citet{jelizarow2010overoptimism}
showed that re-evaluating a published positive result without its optimistic
selection step can remove the advantage entirely. \citet{kapoor2023leakage}
survey how widely the problem persists across applied machine-learning research
and how often it is invisible in the write-up. The cause is that a table of candidate
values reads as a sensitivity analysis whether or not the reported
configuration was picked from it. The distinction that matters is not whether
the parameter grid is shown but whether
the reported configuration was chosen using the labels it is then scored
against. Section~\ref{sec:fusion-protocol} sets out the two protocols we compare
precisely enough for the difference to be checked.
